% Styled after the Magnecko Sense manual/template.tex.
% Compile with LuaLaTeX or XeLaTeX; no external images are required.
\documentclass[10pt,a4paper,oneside]{report}
\usepackage{fontspec}
\setmainfont{DejaVu Sans}
\setsansfont{DejaVu Sans}
\setmonofont[Scale=0.85]{DejaVu Sans Mono}
\usepackage[a4paper,left=20mm,right=19mm,top=25mm,bottom=23mm,
  headheight=28pt,headsep=7mm,footskip=12mm]{geometry}
\usepackage[table]{xcolor}
\usepackage{microtype}
\usepackage{amsmath,amssymb}
\usepackage{array,booktabs,longtable,calc}
\usepackage{graphicx}
\usepackage{fvextra}
\usepackage{enumitem}
\usepackage{titlesec}
\usepackage{fancyhdr,lastpage}
\usepackage{etoolbox}
\usepackage{xurl}
\usepackage{hyperref}
\usepackage{bookmark}
\definecolor{MGSRed}{HTML}{B71930}
\definecolor{MGSInk}{HTML}{1C2630}
\definecolor{MGSMuted}{HTML}{52616D}
\definecolor{MGSRule}{HTML}{C9D0D5}
\definecolor{MGSGray}{HTML}{F0F3F5}
\color{MGSInk}
\newcommand{\ProductName}{Magnecko EPM Controller}
\newcommand{\SourceBaseline}{105e29d}
\newcommand{\DocumentID}{EPM-UG-001}
\newcommand{\DocumentRevision}{A}
\newcommand{\DocumentDate}{October 2026}
\hypersetup{colorlinks=true,linkcolor=MGSRed,urlcolor=MGSRed,
  pdftitle={Magnecko EPM Controller -- User Guide and Reference Manual},
  pdfauthor={Magnecko},bookmarksnumbered=true,bookmarksopen=false}
\urlstyle{tt}
\setlength{\parindent}{0pt}
\setlength{\parskip}{5pt}
\setlength{\emergencystretch}{2em}
\setlength{\tabcolsep}{4pt}
\renewcommand{\arraystretch}{1.18}
\setlength{\LTpre}{8pt}
\setlength{\LTpost}{8pt}
\setcounter{tocdepth}{1}
\setcounter{secnumdepth}{2}
\makeatletter
\renewcommand*\l@section{\@dottedtocline{1}{1.5em}{3.6em}}
\patchcmd{\l@chapter}{1.5em}{2.4em}{}{\PackageError{epmmanual}{Could not size chapter numbers}{Check the table of contents layout.}}
\makeatother
\setlist{nosep,leftmargin=1.5em,topsep=4pt,parsep=2pt}
\setlist[itemize]{label=\textcolor{MGSRed}{\textbullet}}
\providecommand{\tightlist}{\setlength{\itemsep}{1pt}\setlength{\parskip}{1pt}}
\DefineVerbatimEnvironment{verbatim}{Verbatim}{breaklines=true,breakanywhere=true,
  fontsize=\small,baselinestretch=1,frame=single,rulecolor=\color{MGSRule},
  framesep=6pt,formatcom=\color{MGSInk}}
\AtBeginEnvironment{longtable}{\small}
\titleformat{\chapter}[display]{\bfseries\color{MGSRed}}
  {\large\chaptertitlename\ \thechapter}{5pt}{\LARGE}
\titlespacing*{\chapter}{0pt}{0pt}{17pt}
\titleformat{\section}{\large\bfseries\color{MGSRed}}{\thesection}{0.6em}{}
\titleformat{\subsection}{\normalsize\bfseries}{\thesubsection}{0.6em}{}
\titlespacing*{\section}{0pt}{12pt}{5pt}
\titlespacing*{\subsection}{0pt}{9pt}{4pt}
\renewcommand{\chaptermark}[1]{\markboth{\thechapter\quad #1}{}}
\pagestyle{fancy}
\fancyhf{}
\fancyhead[L]{\textcolor{MGSRed}{\large\bfseries EPM}\quad\small\ProductName}
\fancyhead[R]{\scriptsize\DocumentID\ Rev.\ \DocumentRevision\\Source: \SourceBaseline}
\fancyfoot[L]{\scriptsize\color{MGSMuted}\nouppercase{\leftmark}}
\fancyfoot[R]{\small\bfseries\thepage}
\renewcommand{\headrulewidth}{0.5pt}
\renewcommand{\footrulewidth}{0.3pt}
\fancypagestyle{plain}{\fancyhf{}
  \fancyhead[L]{\textcolor{MGSRed}{\large\bfseries EPM}\quad\small\ProductName}
  \fancyhead[R]{\scriptsize\DocumentID\ Rev.\ \DocumentRevision\\Source: \SourceBaseline}
  \fancyfoot[L]{\scriptsize\color{MGSMuted}\DocumentDate\quad |\quad PRELIMINARY}
  \fancyfoot[R]{\small\bfseries\thepage}}
\newcommand{\code}[1]{\texttt{\detokenize{#1}}}
\newcolumntype{L}[1]{>{\raggedright\arraybackslash}p{#1}}

\begin{document}
\raggedright
\hypersetup{pageanchor=false}
\begin{titlepage}
\thispagestyle{empty}
{\color{MGSRed}\Large\bfseries magnecko}\par
\vspace{22mm}
{\color{MGSRed}\small\bfseries MAGNET CONTROL + FOOT SENSING}\par
\vspace{4mm}
{\fontsize{32}{38}\selectfont\bfseries\ProductName}\par
\vspace{6mm}
{\fontsize{22}{29}\selectfont User Guide and\\Reference Manual\par}
\vspace{10mm}
{\large Development, integration, diagnostics, and validation}\par
\vspace{9mm}
\textcolor{MGSRed}{\rule{\linewidth}{1.2pt}}\par
\vspace{8mm}
\begin{tabular}{@{}ll}
Document & \DocumentID\quad Revision \DocumentRevision \\
Firmware baseline & \texttt{\SourceBaseline} \\
Publication & \DocumentDate \\
Status & Preliminary engineering reference
\end{tabular}
\vfill
{\small This manual covers build and programming workflows, EtherCAT
configuration, magnet control, sensor telemetry, diagnostics, and recovery
for the Magnecko EPM controller. Detailed developer documentation is
maintained alongside the firmware in the project wiki.\par}
\vspace{6mm}
{\small\color{MGSMuted}Firmware repository: \texttt{magnecko-epm-controller}\\
Source paths are relative to the repository root. Chapter links are internal;
source-derived behavior does not record completed hardware tests.\par}
\end{titlepage}
\hypersetup{pageanchor=true}
\pagenumbering{roman}
\tableofcontents
\clearpage
\pagenumbering{arabic}
\chapter{System overview}

This guide explains how to prepare, build, commission, and operate the
STM32H725 firmware for a Magnecko electro-permanent magnet (EPM) foot.
It describes the behavior present in the source, including limitations
that affect commissioning and recovery. It is intended for an operator
or integrator with access to the board design and an EtherCAT master.

\textbf{Commissioning status.} The checked-in EtherCAT XML describes an
older actuator interface and does not match the compiled EPM dictionary.
A clean firmware build with ARM GCC 14.3.1 also fails on an existing SPI
receive-buffer pointer mismatch. Resolve those issues before commissioning
a newly built image or master configuration. This guide does not claim a
successful hardware validation.


The controller receives control words and magnet commands through an
external EtherCAT slave controller connected over SPI. Two timer outputs
generate magnetization and demagnetization pulses. A finite state machine
(FSM) gates network pulse requests, and a physical button provides a
separate demagnetization path.

The firmware publishes magnetic holding-force estimates, temperature,
IMU reports, an 8 by 8 distance frame and signal-to-noise values,
four inductive sensor readings, and four Hall sensor readings. A capacitor
charger interface is present, but the current source selects battery mode
and skips charger initialization.

The project name in the build system remains
\code{MagneckoFoot_Header}. The application runs without an RTOS:
the control and communication work runs at a nominal 1\,kHz in TIM2,
IMU service at 400\,Hz in TIM4, and Hall measurements in the foreground
loop. Telemetry may contain the latest stored sample rather than a new
measurement on every EtherCAT cycle.

\section{What this repository supplies}

The repository includes application firmware, STM32 HAL/CMSIS sources,
the CubeMX project, an SOES slave-stack submodule, build files, and
EtherCAT description artifacts. Detailed developer documentation is in
\path{wiki/Home.md}.

Board connector drawings, approved supply voltages, current limits,
mechanical load limits, an EtherCAT master program, and a verified
external ESC EEPROM provisioning procedure are not supplied. Obtain
the applicable board design information before connecting hardware.

\section{Physical state and software state}

Support the mechanism independently before testing release. A
demagnetization pulse, including one requested with the physical button,
can release the magnetic attachment.

The reported magnet status is a software record of the last accepted
pulse direction. It does not confirm magnetic attachment, pulse energy,
or actual holding force. It initializes to false after MCU reset
regardless of the physical magnetic state. Treat the force estimate as
an empirical measurement requiring validation for the individual foot
and surface.

\chapter{Prepare the development environment}

\section{Tools and connections}

Use Git, CMake 3.22 or newer, Ninja, and an ARM embedded GCC toolchain.
Place \code{arm-none-eabi-gcc}, \code{arm-none-eabi-g++}, and binutils
on \code{PATH}.

Flashing requires STM32CubeProgrammer and an ST-Link
connected through the board's SWD interface. Operation requires a
compatible EtherCAT master and a corrected device/process configuration.

The toolchain targets Cortex-M7 with hard-float \code{fpv5-d16}. C is
configured as C11. C++ source uses C++17 constructs, although the build
does not explicitly select a C++ standard. Use a compiler whose default
supports those constructs. C++ exceptions and RTTI are disabled.

The CubeMX project identifies \code{STM32H725VGH6}; the checked-in
VS Code debugger configuration names \code{STM32H725RGVx}. Verify the
physical device and correct the debugger selection before relying on it.

\section{Initialize the slave-stack dependency}

Run from the repository root:

\begin{verbatim}
git submodule update --init --recursive
git submodule status
\end{verbatim}

The recorded SOES URL uses SSH. If SSH access is unavailable, a local
override can select HTTPS without editing the repository's submodule file:

\begin{verbatim}
git config submodule.Libs/SOES_Wrapper/SOES.url https://github.com/OpenEtherCATsociety/SOES.git
git submodule update --init --recursive
\end{verbatim}

Keep the recorded submodule revision unless intentionally reviewing an
upgrade.

\chapter{Build and flash the firmware}
\label{sec:build}

\section{Select the foot identity and build}

Use the CMake presets from the repository root:

\begin{verbatim}
cmake --list-presets
cmake --preset Debug -DEPM_NUMBER=3
cmake --build --preset Debug --parallel
\end{verbatim}

Replace \code{3} with the identity expected by the host for the target
foot. CMake defaults to 3, while the VS Code configuration task offers
1--4 and defaults to 1. Explicit selection avoids unintentionally
programming the wrong identity.

The expected application image is
\path{build/Debug/MagneckoFoot_Header.elf}. The build also produces a map
file and \path{build/Debug/generated/version.h}. The version header records
the eight-character Git hash, a dirty flag, the Git version string, and
the EPM number. Uncommitted documentation also marks the version dirty.

Other configure/build preset pairs are \code{RelWithDebInfo},
\code{Release}, and \code{MinSizeRel}. For example:

\begin{verbatim}
cmake --preset Release -DEPM_NUMBER=3
cmake --build --preset Release --parallel
\end{verbatim}

\textbf{Existing build issue.} A clean Debug configure with ARM GCC
14.3.1 succeeded during documentation checks, but compilation failed in
\path{Libs/SOES_Wrapper/hal/esc_hw.c}, at the SPI call in
\code{esc_address()}. The receive argument \code{&ESCvar.ALevent} does
not have the byte-pointer type HAL expects. This requires a reviewed
firmware correction; the documentation leaves source behavior unchanged.

\section{Program a successfully built image}

Confirm the target identity, board variant, and power arrangement. The
following commands mirror the repository's flash task:

\begin{verbatim}
STM32_Programmer_CLI --list
STM32_Programmer_CLI --connect port=swd --download build/Debug/MagneckoFoot_Header.elf -hardRst -rst --start
\end{verbatim}

The ELF includes its load addresses. Programming it updates the MCU;
it does not program the external EtherCAT controller's EEPROM/SII.

VS Code offers a \code{Build + Flash} task that selects the EPM number,
configures, rebuilds, and flashes. Its launch and attach configurations
use the STM32 extension's CubeCLT paths and Cortex-Debug. The application
has SWO utilities but no UART command shell.

\chapter{Configure the EtherCAT master}
\label{sec:ethercat}

\section{Reconcile the device description}

\textbf{Do not assume the current XML matches the firmware.}
The file \path{Libs/SOES_Wrapper/autogenerated/slave.xml} describes
\code{Actuator V1} with position, velocity, torque, and GPIO mappings.
CMake compiles \path{slave_objectlist.c} from the same directory, whose
interface is the EPM layout below. JSON metadata also differs from the
compiled identity values.

Before commissioning, reconcile the master ESI, actual ESC EEPROM/SII,
and compiled C dictionary. The repository does not supply a verified
end-to-end regeneration/provisioning workflow. Use the compiled dictionary
and \path{utypes.h} as references for the firmware interface.

The application polls SOES from its nominal 1\,kHz timer. Start with
a nominal 1\,ms master process cycle compatible with that behavior.
Distributed-clock validation is not implemented; enabling DC may prevent
the transition from PREOP to SAFEOP.

\section{Master outputs: 4 bytes}

The slave's RxPDO \code{0x1600} contains the following packed fields.
Offsets describe the EtherCAT process image, not the C object structure.

\begin{longtable}{L{17mm}L{28mm}L{50mm}L{26mm}}
\toprule
Byte & Object & Field & Width \\
\midrule
\endhead
0 & \code{0x6040:00} & Control word & 16 bits \\
2 & \code{0x6060:00} & Operation mode & 8 bits \\
3 & \code{0x7000:00} & Magnet command & 8 bits \\
\bottomrule
\end{longtable}

Keep operation mode at zero. It has storage in the dictionary, but no
controller mode-selection behavior is implemented.

\section{Master inputs: 264 bytes}

The default TxPDO assignment lists \code{0x1A00} through
\code{0x1A06} in this order:

\begin{longtable}{L{22mm}L{15mm}L{18mm}L{86mm}}
\toprule
PDO & Byte & Bytes & Fields in order \\
\midrule
\endhead
\code{0x1A00} & 0 & 10 & Status word, mode display, magnet status, force estimate, temperature \\
\code{0x1A01} & 10 & 20 & Gyro XYZ, linear acceleration XYZ, quaternion RIJK \\
\code{0x1A02} & 30 & 128 & 64 signed 16-bit distances \\
\code{0x1A03} & 158 & 64 & 64 unsigned 8-bit SNR values \\
\code{0x1A04} & 222 & 16 & Four unsigned 16-bit LDC frequency registers, then four RP registers \\
\code{0x1A05} & 238 & 24 & Four signed 16-bit Hall X values, then four Y, then four Z \\
\code{0x1A06} & 262 & 2 & Unsigned 16-bit capacitor voltage \\
\bottomrule
\end{longtable}

Use EtherCAT little-endian decoding for multi-byte values. Do not copy
the in-memory \code{_Objects} structure as the process image.

General status at byte offsets 0--9 is:

\begin{longtable}{L{15mm}L{29mm}L{35mm}L{63mm}}
\toprule
Byte & Object & Type & Meaning \\
\midrule
\endhead
0 & \code{0x6041:00} & 16-bit word & Application FSM status \\
2 & \code{0x6061:00} & 8-bit & Mode display; not updated by current application logic \\
3 & \code{0x2000:00} & Unsigned 8-bit & Software magnet status, zero or one \\
4 & \code{0x2001:00} & REAL32 & Empirical filtered force estimate, intended newtons \\
8 & \code{0x2040:01} & Signed 16-bit & Temperature in hundredths of a degree Celsius \\
\bottomrule
\end{longtable}

Read the error code separately with SDO \code{0x603F:00}; it is not
included in the default input PDOs.

\chapter{Commission and enable the application}
\label{sec:enable}

\section{Initial checks}

Use the board's approved supplies and wiring. Keep the magnet command
at zero during startup. Check that the master mapping agrees with
Chapter~\ref{sec:ethercat}, that the EPM identity at
\code{0x2100:02} matches the device, and that required sensor values
and temperature are plausible.

The firmware waits indefinitely for \code{EEPROM_LOADED} after
resetting the external ESC. If initialization remains in its configuring
pattern, inspect that input and the ESC reset/power arrangement.

Bring EtherCAT to OP using the master's normal state transitions.
EtherCAT OP is separate from the application FSM's operation-enabled
state. Exchange initialized outputs continuously: control word zero,
mode zero, and magnet command zero.

\section{Control-word sequence}

Apply each control word until the returned status masked with
\code{0x007F} matches the expected state:

\begin{longtable}{L{29mm}L{28mm}L{29mm}L{51mm}}
\toprule
Step & Control word & Masked status & Expected state \\
\midrule
\endhead
Initial & \code{0x0000} & \code{0x0050} & Switch on disabled \\
Shutdown command & \code{0x0006} & \code{0x0031} & Ready to switch on \\
Switch on & \code{0x0007} & \code{0x0033} & Switched on \\
Enable & \code{0x000F} & \code{0x0037} & Operation enabled \\
\bottomrule
\end{longtable}

These are the constants implemented by this firmware. The drive-style
command name ``shutdown'' denotes a state transition; it does not mean
that a demagnetization pulse is sent. Bit 7, \code{0x0080}, is the
warning indication and lies outside the state mask.

Continue cyclic exchange and add host timeouts for each transition.
Inspect warnings rather than treating the state alone as proof that
all sensors work. Wait for operation enabled before issuing a magnet
command: entry to that state clears pending requests.

\chapter{Magnet commands and manual release}
\label{sec:pulses}

\section{Request a pulse}

Keep control word \code{0x000F} and operation mode zero while requesting
a network pulse. Magnet command bit 0 requests magnetization; bit 1
requests demagnetization. Requests use a rising edge on each bit.

\begin{enumerate}
\item Send \code{Magnet_Command = 0x00} and allow at least one firmware
processing cycle.
\item Send \code{0x01} to magnetize or \code{0x02} to demagnetize.
Keep it present long enough for the device's poll to observe it.
\item Return the command to zero and allow it to be processed before
another request.
\item Inspect magnet status and sensor response. Allow more than
100\,ms from the last accepted pulse before another request, plus any
board-specific recharge time.
\end{enumerate}

The default pulse parameter is 5000\,$\mu$s (5\,ms). The pulse function
accepts parameters between 1000 and 10000\,$\mu$s inclusive, rejects a
request during dead time or with an active overtemperature flag, and
updates the software magnet flag when it starts the pulse timer.

The lockout is approximately 100\,ms. A blocked request is dropped,
not queued. Retry with a fresh zero-to-one edge after the lockout and
recharge interval. Leaving a command bit high does not repeat pulses.
Do not request both directions together: raising both bits from zero
cancels both requests in the input callback.

Writing the pulse-time SDO does not change the pulse width. The active
width comes from \code{MAGNETIZATION_TIME} in the source.

\section{Conceptual cyclic sequence}

This is pseudocode, not a supplied master API. Each exchange represents
continued cyclic process service, with a timeout and fault handling in
the host:

\begin{verbatim}
exchange(control=0x0000, mode=0, magnet=0)
wait_until((status & 0x007F) == 0x0050)
exchange(control=0x0006, mode=0, magnet=0)
wait_until((status & 0x007F) == 0x0031)
exchange(control=0x0007, mode=0, magnet=0)
wait_until((status & 0x007F) == 0x0033)
exchange(control=0x000F, mode=0, magnet=0)
wait_until((status & 0x007F) == 0x0037)

exchange(control=0x000F, mode=0, magnet=0x01)
exchange(control=0x000F, mode=0, magnet=0x00)
wait_beyond_lockout_and_board_recharge_interval()
exchange(control=0x000F, mode=0, magnet=0x02)
exchange(control=0x000F, mode=0, magnet=0x00)
\end{verbatim}

\section{Physical button}

Pressing the active-low button starts a timer. Releasing it stops the
timer. Holding it for about one second directly attempts a
demagnetization pulse. This path bypasses the application's
operation-enabled state check but still respects pulse limits, dead
time, and the overtemperature flag.

The button is not serviced while the MCU is halted. A blocked button
request is not retried automatically; release and press again if a
retry is appropriate. Do not treat the button as independent of
firmware execution or pulse-energy availability.

\chapter{Read and interpret telemetry}
\label{sec:telemetry}

\section{Temperature}

Divide \code{0x2040:01} by 100 to obtain degrees Celsius; for example,
\code{2534} represents 25.34\,$^{\circ}$C. The active sensor model is
PT1000. Disconnected/out-of-range detection can leave this exported
value unchanged, so inspect the error code as well.

\section{IMU and distance data}

The ten IMU fields at \code{0x6090:01}--\code{0x6090:0A} are raw
signed 16-bit report values: gyro XYZ, linear acceleration XYZ, then
quaternion R/I/J/K. Apply the report-specific fixed-point conversion
in the host. Firmware does not convert these to floating-point SI
values or provide scale metadata.

ToF pixels use row-major order: pixel \code{i} corresponds to
row \code{i / 8}, column \code{i mod 8}, and subindex \code{i + 1}.
Distances are at \code{0x60A0:01}--\code{0x60A0:40}, with SNR at
the matching subindex of \code{0x60A1}. The driver stores unsigned
16-bit distances while the dictionary exports signed 16-bit values.
Confirm sensor units, invalid codes, and physical orientation before
interpreting them.

\section{Inductive and Hall values}

The four frequency/inductance registers are
\code{0x60B0:01}--\code{0x60B0:04}; the four RP registers are
\code{0x60B0:05}--\code{0x60B0:08}. They are raw unsigned 16-bit values,
not converted frequencies in hertz.

Hall data at \code{0x60C0} groups all four X channels, then all four
Y channels, then all four Z channels. PDO values are raw signed
registers. Internally the configured driver converts each axis with
\code{raw / 65536 * 2 * 250}, computes magnitudes, and documents them
as mT. Validate that conversion for the fitted device and calibration.

\section{Force estimate and slip indication}

The Hall model averages four magnitudes, estimates an air gap, and
applies empirical force coefficients from \path{TMAG5273.h}. The
controller multiplies the result by its software magnet flag, keeps
a moving average of up to 20 samples, and publishes a REAL32 value
intended to represent newtons.

Each LDC compares fast and slow filters of its RP reading. Three or
more channels indicating slip force the published estimate to zero.
This does not itself trigger demagnetization or change the FSM.
Because the moving average is not cleared when demagnetizing, old
nonzero samples may briefly remain after magnet status becomes zero.

The active model uses compiled coefficients, not the EEPROM parameter
array or SDO calibration storage. Validate the estimate against the
specific foot and surface before making load decisions.

\section{Charger voltage and freshness}

Capacitor voltage at \code{0x60D0:00} is in 10\,mV counts when a
charger has initialized and communicates successfully. The default
\code{USE_BATTERIES} definition skips that initialization, leaving the
value at zero. Zero in this configuration is not a measured
discharged condition.

Sensor PDOs have no timestamps or per-sample validity bits. Failed
or stalled acquisition may leave old or zero values. Observe both
sensor behavior and diagnostics when checking a newly connected board.

\chapter{Disable and shut down}

For a normal release, independently support the mechanism, request
demagnetization while operation is enabled, return the magnet command
to zero, and verify physical release. Then write control word
\code{0x0000} and wait for masked status \code{0x0050}. The disabled
state disables charger command storage and asserts the discharge GPIO.
Follow the board's physical discharge and power-removal procedure.
Firmware does not confirm discharge completion.

Control word \code{0x0007} disables operation into switched-on;
\code{0x0006} returns toward ready; \code{0x0002} requests quick stop,
with masked status \code{0x0017}. The quick-stop action is empty.
Those state changes do not issue a demagnetization pulse.

The SOES software watchdog has a nominal 100-call interval, about
100\,ms at the intended polling rate. It can move EtherCAT to an
error state, but no application safe-output hook is installed.
Loss of EtherCAT, quick stop, disable, or MCU reset must not be
assumed to cause magnetic release or an application discharge action.

\chapter{Warnings, faults, and recovery}
\label{sec:faults}

\section{Temperature behavior}

For a valid PT1000 reading, at least 65\,$^{\circ}$C sets an
overtemperature warning; at least 80\,$^{\circ}$C requests a fatal
fault reaction. At least 240\,$^{\circ}$C is instead classified as
disconnected/out of range and raises a warning. The source has an
intended clear branch at 50\,$^{\circ}$C, but also clears the flag in
the 65--80\,$^{\circ}$C warning branch. Do not assume reliable
80/50\,$^{\circ}$C hysteresis from the current implementation.

\section{Fatal handler}

The fault-reaction action calls \code{Error_Handler()}. It stops pulse
outputs, asserts discharge, disables the charger command, attempts a
charger exchange, disables interrupts, and halts indefinitely. The
last fault status may not reach the master before communication stops.

Correct the cause and reset or power-cycle the controller. Although
the generic FSM defines fault status \code{0x0018} and fault-reset
command \code{0x0080}, this fatal controller path does not transition
to a serviceable fault state. A control-word reset is not a recovery
mechanism for the halted firmware.

\section{Error-code reference}

Read SDO \code{0x603F:00}. The background callback rebuilds the value
on every invocation, and later checks can replace earlier warnings;
it is not an error history.

\begin{longtable}{L{34mm}L{110mm}}
\toprule
Code & Meaning \\
\midrule
\endhead
\code{0x0000} & No currently reported error \\
\code{0x2220} & General overcurrent enum; not currently raised by the controller background logic \\
\code{0x4300} & Overtemperature warning or fault, depending on threshold \\
\code{0xFF01}--\code{0xFF0F} & Combined sensor initialization failures: low-nibble bits IMU=1, LDC=2, ToF=4, Hall=8 \\
\code{0xFF10} & Timer initialization failure; sets fault bit, without directly entering fatal reaction \\
\code{0xFF11} & Invalid EEPROM parameters; validation is currently a stub \\
\code{0xFF12} & Charger overcurrent \\
\code{0xFF13} & Charger overvoltage \\
\code{0xFF14} & Charger watchdog fault \\
\code{0xFF15} & Charger not responding during initialization \\
\code{0xFF16} & Temperature sensor disconnected or out of range \\
\code{0xFF17} & Gate-driver fault while enabled \\
\code{0xFF18} & Invalid simultaneous requests seen by the operation action \\
\bottomrule
\end{longtable}

For example, \code{0xFF05} indicates IMU and ToF initialization
failure. Raising both magnet bits from zero normally cancels requests
in the input callback, so it need not produce \code{0xFF18}.

\section{LED indications}

Initialization selects an alternating blue/green pattern, followed by
cyan. Entry to switched-on selects all channels on; operation-enabled
selects solid green. A fatal error selects fast red blinking.
Sensor-specific animations exist in the driver, but the current
initialization-warning path does not select them. Green is therefore
not proof that every sensor initialized correctly.

\chapter{Configuration and troubleshooting}

\section{Settings that require source changes}

\code{EPM_NUMBER} is a CMake cache setting. Battery mode,
\code{MAGNETIZATION_TIME}, \code{DEAD_TIME}, force-model coefficients,
temperature type, and temperature thresholds are source-level settings.
Passing \code{-DUSE_BATTERIES=OFF} to CMake does not remove the
unconditional header definition.

The following writable SDOs currently provide storage with limited
effects:

\begin{longtable}{L{32mm}L{112mm}}
\toprule
Object & Current behavior \\
\midrule
\endhead
\code{0x2100:01} & Pulse-time storage is not used by pulse generation. \\
\code{0x2100:02} & EPM identity storage is initialized from the build; writes are not persisted. \\
\code{0x2101:01}--\code{0x2101:FA} & 250 REAL32 storage elements; beginning is repurposed for Git version bytes at startup. No active model update or persistence hook is installed. \\
\bottomrule
\end{longtable}

Ten floats are read from the I2C parameter EEPROM at startup, but
validation is unconditional and those values are not used by the
current force formula. Defining the compiler symbol \code{EEPROM_WRITE}
causes an embedded image to be written on every boot of that firmware;
reserve it for deliberate provisioning after reviewing the image.

\section{Symptom checklist}

\begin{longtable}{L{51mm}L{93mm}}
\toprule
Symptom & Check \\
\midrule
\endhead
Missing SOES sources & Initialize the Git submodule and verify SSH access or a local HTTPS override. \\
SPI compile error with GCC 14 & Existing receive-buffer pointer mismatch in \code{esc_address()}; see Chapter~\ref{sec:build}. \\
Configuration pattern never finishes & Inspect \code{EEPROM_LOADED}, ESC reset, and power; startup has no readiness timeout. \\
Master mapping does not match & Reconcile XML, compiled dictionary, and actual SII before commissioning. \\
OP reached but no magnet pulse & Check application state, command edge, lockout, operation-entry timing, and temperature flags. \\
Repeated request ignored & Send zero and allow it to be processed before a new rising edge. \\
Force estimate zero & Inspect magnet flag, Hall data/initialization, and three-of-four LDC slip votes. \\
Capacitor voltage zero & Default battery mode does not initialize charger telemetry. \\
Temperature or sensor value stale & Check error code and driver/interrupt state; PDOs have no freshness flags. \\
Red blinking, no communication & Fatal halt; correct the cause and reset. The manual button cannot run while halted. \\
\bottomrule
\end{longtable}

\section{Implementation issues for integrators}

The ESC read helper has a 128-byte termination buffer although
mailboxes are configured for 512 bytes, and the read-length assertion
is disabled. Review that path before using large mailbox transfers.
The LDC read function supplies a stack-local transmit buffer to a
nonblocking SPI transfer, which also merits review when investigating
corrupt readings.

Custom \code{NOLOAD} memory sections, some filter members without
explicit initialization, shared foreground/interrupt data, blocking
operations, and dynamic allocation in the control interrupt need
hardware validation. Nominal timer rates alone do not establish
execution-time guarantees. These observations are documented without
changing the firmware.

\appendix
\chapter{Source references}

\begin{longtable}{L{66mm}L{78mm}}
\toprule
Path relative to repository root & Purpose \\
\midrule
\endhead
\path{Core/Src/main.cpp} & Entry point, command edge detection, HAL callbacks, fatal handler \\
\path{Core/Inc/main.h} & GPIO and timer aliases \\
\path{Core/Src/tim.c} & Timer prescalers and periods \\
\path{Libs/Controller/FootController.cpp} & Initialization, state actions, pulse generation, telemetry, temperature checks \\
\path{Libs/Controller/FootController.h} & Timing constants, battery selection, and SOES configuration \\
\path{Libs/FSM/FSMTypes.hpp} & Control-word masks and status constants \\
\path{Libs/FSM/FSM.hpp} & Error-code enumeration \\
\path{Libs/SOES_Wrapper/autogenerated/slave_objectlist.c} & Compiled PDO/SDO dictionary \\
\path{Libs/SOES_Wrapper/autogenerated/utypes.h} & Object storage types \\
\path{Libs/SOES_Wrapper/hal/esc_hw.c} & SPI ESC access \\
\path{Libs/Sensors/TMAG5273.cpp} & Hall conversion and empirical force model \\
\path{Libs/Sensors/LDC1101.cpp} & Inductive acquisition and slip filters \\
\path{CMakeLists.txt} & Application sources and EPM selection \\
\path{CMakePresets.json} & Configure/build presets \\
\path{wiki/Home.md} & Full codebase wiki and developer workflow \\
\bottomrule
\end{longtable}

\end{document}
